We introduced a dataset-regime benchmark for empirical quantum reservoir advantage against echo state networks. The main result is deliberately two-sided. ESNs remain the stronger default: QRC does not win on average under the globally frozen protocol. At the same time, QRC wins are reproducible, structured, and predictable from measured dataset properties. The strongest and most likely QRC-usefulness regime in the current evidence is slow, stateful forecasting: drifting trends, regime changes, persistent low-frequency structure, long memory, and temporally correlated fluctuations.

The benchmark therefore turns QRC selection into a measurable triage problem. Practitioners need not ask whether QRC is universally better. They can profile a dataset, compare it with the atlas, and decide whether quantum-reservoir evaluation is worth the cost relative to a strong ESN baseline. This is the practical meaning of quantum reservoir advantage in this paper, and the released code exposes it as a CSV-based triage tool for new univariate forecasting datasets.

The next scientific step is mechanistic. The atlas identifies where QRC helps; it does not yet prove why. Follow-up work should measure reservoir-state memory kernels, latent-state recoverability, phase/coupling/dissipation ablations, and matched long-memory ESN controls in the high-usefulness pockets identified here. Those studies can test whether the observed slow-structure regimes arise from genuinely quantum feature dynamics or from a broader reservoir-memory design principle.

The strongest claim supported now is therefore not universal quantum advantage. It is a regime-map claim: empirical quantum reservoir advantage is selective, measurable, and predictable under a fair frozen comparison against an ESN. That claim is useful even when QRC loses on average, because it gives the field a concrete way to decide where QRC evidence should be concentrated next.
